\documentclass[11pt]{letter}

\usepackage[margin=1in]{geometry}
\usepackage[T1]{fontenc}
\usepackage[utf8]{inputenc}
\usepackage{newtxtext}
\usepackage{parskip}

\signature{Finn Beruldsen}
\address{Finn Beruldsen\\
Center for Nuclear Receptors and Cell Signaling\\
Department of Biology and Biochemistry\\
University of Houston, Texas, USA}
\date{June 1, 2026}

\begin{document}

\begin{letter}{Dear Editor-in-Chief,}

\opening{}

Please consider the manuscript, ``Existing Multi-Material FFF Printers Are Sufficient For (Near) Arbitrary Coloration,'' as a Research Paper for publication in \textit{Computers \& Graphics}.

LayerLoom is a computational fabrication workflow that uses ordinary 3MF geometry to encode layer-wise material schedules for multi-material fused filament fabrication. Instead of assigning each displayed color to a separate filament, LayerLoom converts colored regions into interleaved Z-bands assigned to the physical filaments already loaded in the printer. The resulting color-generating schedule travels with the model itself and can be processed by standard multi-material slicers as ordinary filament-assigned geometry.

The manuscript contributes a geometry-level representation for color scheduling, an algorithmic workflow for converting colored 3MF or GLB models into slicer-ready woven 3MF assemblies, and an experimental evaluation of the resulting printed color states. The validation includes printed gamut samples, calibrated CIE L*a*b*/CIEDE2000 image-based measurements, application-oriented models, slicer compatibility tests, and print-time/material-use benchmarks.

We believe the manuscript is well suited to \textit{Computers \& Graphics} because it frames color expansion in FFF as a graphics and computational-fabrication problem: how to encode appearance and material scheduling into a portable model representation rather than into custom hardware, custom G-code, or slicer-specific color-mixing tools. This is closely aligned with prior work in 3D printing, appearance reproduction, geometry processing, and fabrication-aware graphics workflows.

Reproducibility is supported through the manuscript, Supplementary Information, and an anonymized review archive. The software release, example 3MF files, analysis scripts, benchmark data, and figure-generation workflows will be archived with a persistent DOI before publication.

This manuscript has not been published and is not under consideration elsewhere.

Thank you for your consideration.

\closing{Sincerely,}

\end{letter}

\end{document}
